% 【综合论述】所属：第3章（由 chapters/revised/03-justification-context.tex \input 引入）
\minitopic{“有理由”至少包含两个问题}主体可能拥有支持$p$的证据，却因愿望、偏见或迷信而相信$p$；此时命题得到确证，具体信念却没有适当基于证据。反过来，一个过程偶然从某条证据开始，也不表示主体能响应证据变化。基于关系因此需要超出时间先后：理由必须在信念形成或维持中发挥认识作用，并使主体在理由削弱、反向或被替换时以相称方式调整。因果、反事实和规范三类理论各抓住一面，任何单一标准都要面对偏离因果、过度心理化或难以操作的问题。

\minitopic{回溯理论在回答不同疑问}基础主义问理由链怎样获得非推论起点；融贯主义问一个信念如何由整个系统的解释联系支持；无限主义则拒绝把可继续给出理由视为缺陷。它们常被当作互斥标签，但现实认识可以包含多种结构：经验提供可错的初始约束，背景信念决定经验怎样分类，推理又把局部支持整合进更大网络。争点不是能否在一幅图中同时画出基础和融贯，而是哪种结构在提供最终规范地位、如何防止一个内部一致却脱离世界的系统自我支持。

\minitopic{内外之争应按认识地位拆分}恶魔受害者与正常主体具有相同内部经验和推理，却处于完全不同的可靠环境。若问题是主体是否可责、是否从自己的角度理性，两者似乎相同；若问题是信念是否由可靠方式连接事实，两者明显不同。把所有评价都叫作“确证”，争论会变成术语对撞。更精确的理论可承认主观合理性与知识资格分层：内部证据解释主体为何有理由相信，外部可靠联系解释信念为何成功。真正难点是两层怎样互动，而不是简单宣布其中一层不存在。

\minitopic{风险不会神秘地改变证据}高风险案例常让人要求更多复核，但至少有四种解释。风险可能促使主体取得新证据；可能使一个原本可忽略的替代情景成为高阶击败者；可能只改变行动的最优方案；也可能真正提高知识门槛。只有最后一种支持实践侵入。案例若说“证据完全相同”，却又让高风险主体想到新的失败可能，就已经暗中改变了主体信息。严谨比较需要分别固定证据、注意、说话语境和行动代价，观察知识、断言、信念理性与行动许可是否一起变化。

\begin{argumentbox}
一个完整的理由档案应记录：主体在何时拥有何种证据；哪项证据实际参与信念形成；形成过程是否能响应证据变化；当时有哪些一阶与高阶击败者；主体为何停止调查；信念将用于什么断言或行动。缺少任一栏，都容易把不同认识评价混成一句“理由够不够”。
\end{argumentbox}

\minitopic{击败是动态关系而非红色印章}得知钟已停摆直接削弱“现在九点”的依据；得知消息来自竞争者可能只削弱来源可靠性；后来发现竞争者说谎，又可能恢复原证据地位。击败者本身也可能被击败，认识状态因而随时间变化。误导性击败者尤其重要：主体在当时合理降级信念，即使事后发现新信息为假。只按最终真值评价会产生后见偏差，把谨慎修正误作无知，把幸运坚持误作坚定。

\minitopic{公共理由不要求所有因素都可内省}科学报告和司法判决需要公开理由，但公开性可以由制度实现，而不要求每位主体反思掌握全部可靠条件。校准记录、数据出处、预注册和复核流程使外部可靠性部分变得可检查；同时，主体仍须能够理解核心证据为何相关并响应反证。这样一种分层责任观避免两个极端：把不可内省的可靠性视为毫无认识意义，或把系统可靠性当作个人无需思考的免责凭证。

\minitopic{阶段性结论}理由理论的共同任务是解释支持怎样获得、怎样进入信念、怎样被新信息撤销以及怎样许可言说和行动。基础、融贯、内部可及与外部可靠分别提供结构资源，语境和风险则检验这些资源在使用中的边界。成熟立场应明确自己评价的是命题、具体信念、主体理性还是知识，并让不同层次能够在同一时间线上接受检验。
